\section{Normalization of the transfer weight}
\label{sec:supp-weight-normalization}

The population criterion uses
\[
v_{M,h}(W)
=
\frac{h^dK_h^2(X-x_0)}{\mu_h^2}
\left\{\frac1{\pi_M(W)}-1\right\}.
\]
The coefficient and transfer sign are invariant to a positive normalization that is common to every observation and threshold.

\begin{lemma}[Weight-normalization invariance]
\label{lem:weight-normalization}
Let \(c_{M,h}>0\) and replace \(v_{M,h}\) by \(v_{M,h}^{c}=c_{M,h}v_{M,h}\). Then
\[
J_{M,h}^{c}(\lambda)=c_{M,h}J_{M,h}(\lambda),
\qquad
G_{M,h}^{c}=c_{M,h}G_{M,h},
\qquad
H_{M,h}^{c}=c_{M,h}H_{M,h}.
\]
Consequently,
\[
\Pi_{[0,1]}(G_{M,h}^{c}/H_{M,h}^{c})
=
\lambda_h^{\mathrm{or}},
\]
and
\[
\operatorname{sign}(2G_{M,h}^{c}-H_{M,h}^{c})
=
\operatorname{sign}(2G_{M,h}-H_{M,h}).
\]
The same identities hold for the selection-fold empirical criterion under a positive fold-specific constant.
\end{lemma}

\begin{proof}
Each criterion quantity is linear in the weight. The ratio and sign conclusions follow because \(c_{M,h}>0\).
\end{proof}

Taking
\[
c_{M,h}=\mu_h^2/h^{2d}
\]
produces the equivalent population weight
\[
h^{-d}K_h^2(X-x_0)
\left\{\frac1{\pi_M(W)}-1\right\}.
\]
On a selection fold, the implementation uses
\[
\widehat v_h(W)
=
\frac{h^dK_h^2(X-x_0)}{\widehat\mu_h^{\,2}}
\left\{\frac1{\pi_M(W)}-1\right\}.
\]
This equals the empirical version of the equivalent \(h^{-d}\) weight above
multiplied by the positive fold-specific constant
\(h^{2d}/\widehat\mu_h^{\,2}\). The constant changes the numerical scale of
\(\widehat G_h\) and \(\widehat H_h\), but not their ratio or the transfer sign.
